\bookchapter{Chapter 7: Comprehensive Analysis of Cutting-edge Features}

\chapterintro{Supported models:}

\booksection{Apple}

: iPhone 14 series and later support satellite emergency SOS. iPhone 16 series further supports satellite iMessage and SMS.

\booksubsection{Supported models:}

: Mate 50 series started supporting BeiDou satellite messages (text only); Mate 60 Pro series further supports Tiantong satellite calls (can make phone calls); Mate 70 series has further enhanced satellite communication capabilities.

\booksubsection{Things to note:}

\begin{itemize}
  \item \textbf{Apple}: Starting from the iPhone 14 series, satellite emergency SOS is supported. The iPhone 17 series further supports sending iMessages and text messages via satellite, with expanded coverage in more regions.
  \item \textbf{Huawei}Use cases:
  \item \textbf{Outdoor adventure, mountaineering emergencies.}Maritime emergencies.
  \item \textbf{Natural disasters causing ground communication outages.}How to use:
\end{itemize}

\booksubsection{Make sure you're in an open area with a view of the sky.}

\begin{itemize}
  \item Turn on your phone's satellite communication feature (usually in the emergency call interface).
  \item Follow the on-screen instructions to aim at the satellite.
  \item Send preset distress messages and your location.
\end{itemize}

\booksubsection{Trivia}

\begin{enumerate}
  \item Satellite communication requires special hardware antennas in the phone—not all phones support it. Additionally, satellite communication usually requires a paid subscription service, but the emergency SOS feature is generally free.
  \item II. Off-Grid Communication: A New Way of Close-Range Communication
  \item "Off-grid communication" refers to technology that allows two phones to communicate directly with each other without mobile networks or Wi-Fi.
  \item Technical principle:
\end{enumerate}

\begin{tcbinfo}{Uses Bluetooth or dedicated radio frequencies to establish point-to-point connections over short distances (usually a few hundred meters to a few kilometers).}
Application scenarios:

\end{tcbinfo}

\booksection{Outdoor activities}

: Communication between teammates during hiking and camping.

\booksubsection{Disaster relief}

: When communication is disrupted after earthquakes and floods.

\booksubsection{Smart home}

\begin{itemize}
  \item \textbf{: When crowds at concerts and sporting events cause poor signal.}Supported platforms:
  \item \textbf{Apple}: iOS 17's "NameDrop" and "Find Friends Offline" features.
  \item \textbf{Huawei}: The Celia Chat app supports off-grid calls.
\end{itemize}

\booksubsection{Xiaomi}

\begin{itemize}
  \item \textbf{Apple}III. Crash Detection: A Lifesaver at Critical Moments
  \item \textbf{Huawei}Working principle:
  \item \textbf{Xiaomi}Pop up a countdown interface, asking if you're okay.
\end{itemize}

\booksection{If you don't respond before the countdown ends, automatically dial emergency numbers.}

Send your location to emergency contacts.

\booksubsection{Supported models:}

Apple

\begin{enumerate}
  \item : iPhone 14 series and later.
  \item Android
  \item : Some high-end models have successively added support.
\end{enumerate}

\booksubsection{Things to note:}

\begin{itemize}
  \item \textbf{Apple}May trigger false alarms (like roller coasters, intense exercise)—just cancel in time.
  \item \textbf{Android}IV. SparkLink Technology: Next-Generation Wireless Connectivity
\end{itemize}

\booksubsection{SparkLink is a new-generation wireless communication technology led by Huawei, aiming to replace Bluetooth and Wi-Fi in certain scenarios. It matches or even exceeds existing solutions on key metrics: air interface rate up to 900Mbps (far exceeding Bluetooth's 50Mbps), latency as low as 20 $\mu$ s (Bluetooth about 15ms), and a single network supports 4096 devices connected simultaneously (Bluetooth usually under 8).}

\begin{itemize}
  \item Technical features:
  \item Ultra-low latency
  \item : End-to-end latency of only 20 microseconds (0.02ms), two orders of magnitude lower than Bluetooth—gaming and music sync are virtually imperceptible.
\end{itemize}

\booksection{High speed}

: Single-link peak rate of 900Mbps, dozens of times that of Bluetooth audio, capable of transmitting lossless audio and even HD video.

\booksubsection{Low power consumption}

\begin{itemize}
  \item \textbf{: Standby power consumption comparable to Bluetooth, but with higher transmission efficiency and lower energy consumption per unit of data.}Wide coverage + high concurrency
  \item \textbf{: Connection distance over 50\% farther than Bluetooth, with a single network supporting thousands of devices online simultaneously.}Application scenarios:
  \item \textbf{Wireless earbuds}: Lower latency, perfect audio-video sync for gaming.
  \item \textbf{Smart watches}: More stable connections.
\end{itemize}

\booksubsection{Smart home}

\begin{itemize}
  \item \textbf{: More devices connected simultaneously.}Automotive connectivity
  \item \textbf{: Smoother phone-to-car connections.}Supported devices:
  \item \textbf{Smart home}Technology Outlook
  \item \textbf{SparkLink technology represents the future direction of wireless communication. As more manufacturers join, we may see more devices supporting SparkLink, enjoying a better wireless experience.}Go Check: Does Your Phone Support These Features
\end{itemize}

\booksubsection{Search for “satellite” in settings to see if there are satellite communication options.}

Search for “crash detection” or “emergency SOS” to check if it's enabled.

\textbf{If your phone is Huawei, see if it supports SparkLink device connections.}\par

Learn how to use these features—they may come in handy at critical moments.
